\section{Introduction}
\label{sec:introduction}

Training LLM agents on failure trajectories is an increasingly popular paradigm. PALADIN \cite{paladin2025} demonstrates that taxonomy-guided error injection---where training data includes systematically corrupted trajectories tagged with error types---improves tool-use agent performance by +57pp. Related approaches using failure-augmented training show gains across tool-use \cite{song2024eto,eef2025,wang2024learning_from_failure} and dialogue tasks \cite{agentHER2025}. A natural question arises: does this transfer to embodied planning, where agents must execute multi-step physical manipulation tasks?

We test this hypothesis on ALFWorld \cite{shridhar2021alfworld}, a text-based embodied planning benchmark. We design TaxEI, a four-operator error taxonomy for ALFWorld (step-order swaps, sub-goal deletions, context truncations, observation masking), and train Qwen3-8B agents under three conditions: clean demonstrations only (Vanilla), clean + taxonomy-tagged errors (TaxEI), and clean + generic-labeled errors (Random). Each condition uses three seeds.

The answer is negative: TaxEI yields no improvement ($-8.6$pp, $p{=}0.054$) and directionally increases training instability. But the contribution of this paper is not merely reporting failure. We identify three boundary conditions that characterize \emph{when} error injection helps:

\begin{enumerate}[nosep,leftmargin=*]
\item \textbf{Variance gradient.} Training instability follows an ordering: Vanilla (SD 3.3pp) $<$ Random (7.0pp) $<$ TaxEI (8.6pp). While all variance tests are NS at $n{=}3$, the directional pattern suggests taxonomy labels amplify seed sensitivity.

\item \textbf{Dose-response threshold.} Both TaxEI and Random share a sharp performance collapse between 10\% and 25\% error injection ($-12.8$pp each), followed by a plateau. Below this threshold, TaxEI has zero effect while Random provides mild augmentation (+6.4pp).

\item \textbf{Per-operator local/global split.} Local errors (observation masking: $+9.1$pp) are learnable; global errors (context truncation, sub-goal deletion: $-8.3$pp each) are harmful. This suggests error injection requires errors that are observable and locally recoverable---properties present in tool-use but largely absent in embodied planning.
\end{enumerate}

These findings provide actionable guidance: before applying error injection to a new domain, practitioners should assess whether the task's error structure is discrete, observable, and locally recoverable.
